% Bagian: intuitionistic-logic
% Bab: tableaux
% Subbagian: rules

\documentclass[../../../include/open-logic-section]{subfiles}

\begin{document}

\olfileid[id]{int}{tab}{rul}

\olsection{Aturan untuk Logika Intuisionistik}

Aturan-aturan untuk konektif $\land$ dan $\lor$ sama seperti aturan bagi
!!{tableau} proposisional bertanda biasa, hanya dengan prefiks yang
ditambahkan. Dalam setiap kasus, aturan yang diterapkan pada !!{formula}
bertanda $\sFmla{S}{!A}[\sigma]$ menghasilkan !!{formula} baru yang juga
berprefiks~$\sigma$. Hal ini seharusnya jelas secara intuitif: misalnya, jika
$!A \land !B$ benar pada (dunia yang dinamai oleh)~$\sigma$, maka $!A$ dan
$!B$ benar pada~$\sigma$ (dan bukan pada dunia lain). Kita menghimpun aturan
untuk $\land$ dan $\lor$ dalam \olref{tab:prop-rules}.

\begin{table}
  \[\def\arraystretch{3}\begin{array}{|c|c|}
    \hline
    \AxiomC{\sFmla{\True}{!A \land !B}[\sigma]}
    \RightLabel{\TRule{\True}{\land}}
    \UnaryInfC{\sFmla{\True}{!A}[\sigma]}
    \noLine
    \UnaryInfC{\sFmla{\True}{!B}[\sigma]}
    \DisplayProof
    &
    \AxiomC{\sFmla{\False}{!A \land !B}[\sigma]}
    \RightLabel{\TRule{\False}{\land}}
    \UnaryInfC{$\sFmla{\False}{!A}[\sigma] \quad \mid \quad
      \sFmla{\False}{!B}[\sigma]$}
    \DisplayProof
    \\[2ex]
    \hline
    \AxiomC{\sFmla{\True}{!A \lor !B}[\sigma]}
    \RightLabel{\TRule{\True}{\lor}}
    \UnaryInfC{$\sFmla{\True}{!A}[\sigma] \quad \mid \quad
      \sFmla{\True}{!B}[\sigma]$}
    \DisplayProof
    &
    \AxiomC{\sFmla{\False}{!A \lor !B}[\sigma]}
    \RightLabel{\TRule{\False}{\lor}}
    \UnaryInfC{\sFmla{\False}{!A}[\sigma]}
    \noLine
    \UnaryInfC{\sFmla{\False}{!B}[\sigma]}
    \DisplayProof
    \\[2ex]
    \hline
  \end{array}\]
  \caption{Aturan !!{tableau} berprefiks untuk $\land$ dan $\lor$}
  \ollabel{tab:prop-rules}
\end{table}

Syarat penutupan serupa dengan syarat bagi !!{tableau} biasa, meskipun kita
mensyaratkan agar bukan hanya !!{formula}, melainkan prefiksnya juga cocok.
Sesungguhnya, kita dapat sedikit lebih longgar: dalam model intuisionistik,
begitu !!{formula} menjadi benar, formula tersebut tetap benar. Karena itu,
mustahil $!A$ benar pada~$\sigma$ tetapi salah pada prefiks
terjangkau~$\sigma.{*}$. Jadi, suatu cabang tertutup jika memuat keduanya,
\[
\sFmla{\True}{!A}[\sigma] \quad\text{dan}\quad
\sFmla{\False}{!A}[\sigma.{*}]
\]
untuk suatu prefiks~$\sigma$ dan !!{formula}~$!A$. Perhatikan bahwa jika
tandanya dibalik, yakni jika cabang tersebut memuat
\[
\sFmla{\False}{!A}[\sigma] \quad\text{dan}\quad
\sFmla{\True}{!A}[\sigma.{*}],
\]
cabang itu tertutup hanya jika $*$ adalah barisan kosong.

Selain itu, suatu cabang tertutup jika memuat
$\sFmla{\True}{\bot}[\sigma]$.

Aturan untuk menetapkan asumsi juga sama seperti aturan bagi !!{tableau}
biasa, kecuali bahwa untuk asumsi kita selalu menggunakan prefiks~$1$.
(Tidak penting prefiks mana yang kita gunakan, asalkan prefiks tersebut sama
untuk semua asumsi.) Jadi, misalnya, kita mengatakan bahwa
\[
!B_1, \dots, !B_n \Proves !A
\]
jika dan hanya jika terdapat tableau tertutup untuk asumsi-asumsi
\[
\sFmla{\True}{!B_1}[1], \dots, \sFmla{\True}{!B_n}[1],
\sFmla{\False}{!A}[1].
\]

Untuk kondisional~$\lif$, aturannya berbeda dari kasus klasik dan modal.
Aturan $\TRule{\True}{\lif}$ memperluas suatu cabang yang memuat
$\sFmla{\True}{!A \lif !B}[\sigma]$ dengan
$\sFmla{\False}{!A}[\sigma.{*}]$ dan
$\sFmla{\True}{!B}[\sigma.{*}]$ pada dua cabang yang berbeda. Aturan ini
hanya dapat diterapkan untuk prefiks~$\sigma.{*}$ yang \emph{sudah} muncul
pada cabang tempat aturan tersebut diterapkan. Mari kita sebut prefiks
semacam itu ``sudah digunakan'' (pada cabang tersebut). (Karena
$\sigma.{*}$ mencakup $\sigma$ sendiri, aturan itu selalu dapat diterapkan
dengan menambahkan formula bertanda berprefiks
$\sFmla{\False}{!A}[\sigma]$ dan $\sFmla{\True}{!B}[\sigma]$ pada cabang
yang terpisah.)

Aturan $\TRule{\False}{\lif}$ memperluas suatu cabang yang memuat
$\sFmla{\False}{!A \lif !B}[\sigma]$ dengan
$\sFmla{\True}{!A}[\sigma.n]$ dan $\sFmla{\False}{!B}[\sigma.n]$ pada
cabang yang sama, dengan $\sigma.n$ sebagai prefiks yang baru pada cabang
tersebut.

Aturan untuk $\lnot$ didefinisikan secara analog (menggunakan definisi
$\lnot !A$ sebagai $!A \lif \lfalse$).

Aturan-aturan tersebut diberikan dalam \olref{tab:rules-lif-lnot}.

\begin{table}
  \[\def\arraystretch{3}\begin{array}{|c|c|}
    \hline
    \AxiomC{\sFmla{\True}{\lnot !A}[\sigma]}
    \RightLabel{\TRule{\True}{\lnot}}
    \UnaryInfC{$\sFmla{\False}{!A}[\sigma.{*}]$}
    \DisplayProof
    &
    \AxiomC{\sFmla{\False}{\lnot !A}[\sigma]}
    \RightLabel{\TRule{\False}{\lnot}}
    \UnaryInfC{\sFmla{\True}{!A}[\sigma.n]}
    \DisplayProof
    \\[1ex]
    \text{$\sigma.{*}$ sudah digunakan} & \text{$\sigma.n$ baru}\\
    \hline
    \AxiomC{\sFmla{\True}{!A \lif !B}[\sigma]}
    \RightLabel{\TRule{\True}{\lif}}
    \UnaryInfC{$\sFmla{\False}{!A}[\sigma.{*}] \quad \mid \quad
      \sFmla{\True}{!B}[\sigma.{*}]$}
    \DisplayProof
    &
    \AxiomC{\sFmla{\False}{!A \lif !B}[\sigma]}
    \RightLabel{\TRule{\False}{\lif}}
    \UnaryInfC{\sFmla{\True}{!A}[\sigma.n]}
    \noLine
    \UnaryInfC{\sFmla{\False}{!B}[\sigma.n]}
    \DisplayProof
    \\[1ex]
    \text{$\sigma.{*}$ sudah digunakan} & \text{$\sigma.n$ baru}\\
    \hline
  \end{array}\]
  \caption{Aturan !!{tableau} berprefiks untuk $\lnot$ dan $\lif$}
  \ollabel{tab:rules-lif-lnot}
\end{table}

\end{document}
